We will now prove that, given any representation-infinite poset $P$, we can compute a brick-continuous family and a generic brick in $\Mod{\incidencealgebra{P}}$ using the extension functors we just defined.

\Paragraphtitlebf{Reduction of $P$.} Let $P$ be any representation-infinite poset. By Lemma \ref{lemma:infinitereduction}, it can be reduced to a critical poset $C$ (i.e.~a poset with minimally representation-infinite incidence algebra). As done before, we may assume that any step of the reduction process removes only one vertex. Let $\card{P}=n$ and $\card{C}=n-m$, so the reduction process involved $m$ points of $P$ through a sequence of subtractions and/or contractions. We obtain a list of posets $P_i$ such that $\card{P_i}=n-i$ for $0\leq i\leq m$ with $P=P_0$ and $C=P_m$ and we denote by $r_j$ the reduction applied at the $j$-th step ($1\leq j\leq m$), that is 
\[r_j=c\colon P_{j-1} \to P_j\]
if we applied an elementary contraction and
\[r_j=s\colon P_j \to P_{j-1}\]
if we applied a subtraction of a vertex.
From this list of posets and reduction maps we can define a composition of the corresponding extension functors as
% https://q.uiver.app/#q=WzAsNSxbMCwwLCJcXE1vZHtcXGluY2lkZW5jZWFsZ2VicmF7XFxLfXtDfX0iXSxbMSwwLCJcXE1vZHtcXGluY2lkZW5jZWFsZ2VicmF7XFxLfXtQX3ttLTF9fX0iXSxbMiwwLCJcXGNkb3RzIl0sWzMsMCwiXFxNb2R7XFxpbmNpZGVuY2VhbGdlYnJhe1xcS317UF8xfX0iXSxbNCwwLCJcXE1vZHtcXGluY2lkZW5jZWFsZ2VicmF7XFxLfXtQfX0iXSxbMCwxLCJFX3tyX219Il0sWzEsMiwiRV97cl97bS0xfX0iXSxbMiwzLCJFX3tyXzJ9Il0sWzMsNCwiRV97cl8xfSJdXQ==
\[\begin{tikzcd}[ampersand replacement=\&,cramped]
	{\Mod{\incidencealgebra{C}}} \& {\Mod{\incidencealgebra{P_{m-1}}}} \& \cdots \& {\Mod{\incidencealgebra{P_1}}} \& {\Mod{\incidencealgebra{P}}}
	\arrow["{E_{r_m}}", from=1-1, to=1-2]
	\arrow["{E_{r_{m-1}}}", from=1-2, to=1-3]
	\arrow["{E_{r_2}}", from=1-3, to=1-4]
	\arrow["{E_{r_1}}", from=1-4, to=1-5]
\end{tikzcd}\]
that we denote by $\widetilde{E}\defeq E_{r_1}E_{r_2}\cdots E_{r_{m-1}}E_{r_m}$. We will denote by $\widetilde{E}_\downarrow \defeq E_{r_1 \downarrow} E_{r_2\downarrow} \cdots E_{r_{m-1} \downarrow} E_{r_m \downarrow}$ its restriction to finite-dimensional modules.

Let $\{B(\lambda)_C\}_{\lambda\in\K}$ be the brick-continuous family and $G_C$ be the generic brick of the critical poset $C$ given by Theorem \ref{thm:genbrickcritical} (see Appendix \ref{appendix} for the list). By the construction of $G_C$ we have that $\dim_\K((G_C)_i)=\infty$ and $\dim_{\K(x)}((G_C)_i) = (y_C)_i$ for all $i\in (\HasseQuiver{C})_0$ since $(G_C)_i=\K(x)^{(y_C)_i}$ with $y_C$ the minimal sincere positive radical vector of the quadratic form $q_C$ (see Theorem \ref{thm:genbrickcritical}).

Then, with the previous setting, we have the following results:

\begin{lemma}\label{lemma:extendinggenbricks}
	$G_P \defeq \widetilde{E}(G_C)= E_{r_1}E_{r_2}\cdots E_{r_{m-1}}E_{r_m}(G_C)$ is a generic brick for $\incidencealgebra{P}$ with $(G_P)_i =\K(x)^{h_i}$ for all $i\in (\HasseQuiver{P})_0$ for some $h_i >0$ and $\End(G_P)\cong \K(x)$.
\end{lemma}
\begin{proof}
	We prove the statement by induction on the number $m$ of reduction steps.
	Let $m=1$, then we just removed a vertex through a contraction or a subtraction and, up to re-indexing, let $P=\{1,\dots, n-1, n\}$ and $C=\{1,\dots, n-1\}$.
	
	$(i)$ If we applied a contraction $c\colon P\to C$, contracting $n-1$ and $n$, we have $\widetilde{E}=E_c$. Now $(G_P)_i = (G_C)_i$ for the vertices $1,\dots, n-1$ and  $(G_P)_n=(G_C)_{n-1}$ in vertex $n$ by definition. So we get that
	\[\dim_\K(G_P)=\dim_\K(G_C) + \dim_\K((G_C)_{n-1})\]
	Since $G_C$ is a generic brick we have (by \cite[Ch.III, Corollary 3.6]{sim:vol1})
	\[\infty=\ell_{\incidencealgebra{C}}(G_C) =\dim_\K(G_C)\]
	Therefore, we conclude that 
	\[\ell_{\incidencealgebra{P}}(G_P) =\dim_\K(G_P)=\dim_\K(G_C) + \dim_\K((G_C)_{n-1})=\infty\]
	Now $E_c$ is fully faithful (by its definition), so $ \End(G_P)=\End(E_c(G_C)) \cong \End(G_C) \cong \K(x)$. Again by \cite[Ch.III, Corollary 3.6]{sim:vol1}, since $G_C$ has finite endolength we also get 
	\[\begin{aligned}
		\ell_{\K(x)\op}(G_P) &=\dim_{\K(x)}(G_P)=\dim_{\K(x)}(G_C) + \dim_{\K(x)}((G_C)_{n-1}) \\
		&=\ell_{\K(x)\op}(G_C) + (y_C)_{n-1} < \infty
	\end{aligned}\]
	Therefore $G_P$ has infinite length but finite endolength, so it is a generic module and by Remark \ref{rem:genmodulek(x)} it is a generic brick. We also conclude that $(G_P)_i=(E_c(G_C))_i =\K(x)^{h_i}$ for all $i\in (\HasseQuiver{P})_0$ for some $h_i >0$ by definition of $E_c$.
	
	$(ii)$ If we applied a subtraction $s\colon C\to P$, removing the vertex $n$, we have $\widetilde{E}=E_s$. Suppose $\widetilde{E} = E_s = E^{\Hom}_s $ (the case with $E^\otimes_s$ is analogous). Similarly to $(i)$, by definition of $E_s$ we have
	
	\[
	\begin{aligned}
		\dim_\K(G_P)&=\dim_\K(G_C) + \dim_\K(\lim_{n<_P j} G(C)_j) \\
		\dim_{\K(x)}(G_P)&=\dim_{\K(x)}(G_C) + \dim_{\K(x)}(\lim_{n<_P j} G(C)_j)
	\end{aligned}
	\]
	
	\noindent and $\lim_{n<_P j} G(C)_j \leq \prod_{n<_P j} (G_C)_j = \prod_{n<_P j} \K(x)^{(y_C)_j}$ (finite product). Moreover $E_s$ is fully faithful by Lemma \ref{lemma:ffsubtraction}. Therefore, with a similar reasoning as in $(i)$, we conclude that $\ell_{\incidencealgebra{P}}(G_P)=\infty$, $\End(G_P) \cong \End(G_C) \cong \K(x)$ and $\ell_{\K(x)\op}(G_P)<\infty$. Thus, $G_P$ has infinite length but finite endolength, so it is a generic module and by Remark \ref{rem:genmodulek(x)} it is a generic brick. By definition of $E_s$ we also conclude that $(G_P)_i=(E_s(G_C))_i =\K(x)^{h_i}$ for all $i\in (\HasseQuiver{P})_0$ for some $h_i >0$.\smallskip
	
	Now suppose the result holds for $m-1$ reduction steps, so $\widetilde{G}= E_{r_2}\cdots E_{r_{m-1}}E_{r_m}(G_C)$ is a generic brick for $\incidencealgebra{P_1}$ with $(\widetilde{G})_i =\K(x)^{\tilde{h}_i}$ for all $i\in (\HasseQuiver{P})_0$ for some $\tilde{h}_i >0$ and $\End(\widetilde{G})\cong \K(x)$. We have that $r_1$ can either be a contraction or a subtraction but in any case, since $\widetilde{G}$ has a similar behavior as $G_C$, we can apply the same reasoning used in $(i)$ and $(ii)$ to conclude that $G_P=\widetilde{E}(G_C)= E_{r_1}(\widetilde{G})$ is a generic brick for $\incidencealgebra{P}$ with $(G_P)_i =\K(x)^{h_i}$ for all $i\in (\HasseQuiver{P})_0$ for some $h_i >0$ and $\End(G_P)\cong \K(x)$. So we conclude.
\end{proof}

\begin{lemma}\label{lemma:extendingbrickcont}
	The family $\widetilde{E}_\downarrow (B(\lambda)_C) = E_{r_1 \downarrow}E_{r_2\downarrow}\cdots E_{r_{m-1} \downarrow}E_{r_m \downarrow} (B(\lambda)_C)$, for $\lambda\in \K$, contains a brick-continuous family for $\incidencealgebra{P}$ having a fixed dimension vector.	
\end{lemma}
\begin{proof}
	As previously done, we prove the statement by induction on the number $m$ of reduction steps.
	Let $m=1$, then we just removed a vertex through a contraction or a subtraction and, up to re-indexing, let $P=\{1,\dots, n-1, n\}$ and $C=\{1,\dots, n-1\}$.
	
	Let $\underline{d} = (d_1, \dots, d_{n-1}) \defeq y_C = \dimvec{B(\lambda)_C}$ and $d \defeq \dim_\K(B(\lambda)_C) = \sum_{i=1}^{n-1} d_i$.	
	By fully faithfulness we have that $\{\widetilde{E}_\downarrow (B(\lambda)_C)\}_{\lambda\in \K}$ is always an infinite family of non-isomorphic bricks.
	
	$(i)$ If we applied a contraction $c\colon P\to C$, contracting $n-1$ and $n$, we have $\widetilde{E}_\downarrow=E_{c\downarrow}$.
	Then $\underline{d}' \defeq \dimvec{E_{c\downarrow}(B(\lambda)_C)}=(\underline{d}, d_{n-1})$ (and $\dim_\K(E_{c\downarrow}(B(\lambda)_C))= d + d_{n-1}$) for any $\lambda\in\K$. So, for $\lambda\in \K$, $B(\lambda)_P\defeq E_{c\downarrow}(B(\lambda)_C)$ gives a brick-continuous family with dimension vector $\underline{d}'$.
	
	$(ii)$ If we applied a subtraction $s\colon C\to P$, removing the vertex $n$, we have $\widetilde{E}_\downarrow=E_{s\downarrow}$. Suppose $\widetilde{E}_\downarrow = E_{s\downarrow} = E^{\Hom}_{s\downarrow}$ (the case with $E^\otimes_{s\downarrow}$ is analogous). 
	
	Then $\dimvec{E_{s\downarrow}(B(\lambda)_C)}=(\underline{d}, \dim_\K(\lim_{n<_P j} (B(\lambda)_C)_j))$. But now we have that
	\[
	0 \leq \dim_\K(\lim_{n<_P j} (B(\lambda)_C)_j) \leq D_L \defeq \dim_\K(\prod_{n<_P j} (B(\lambda)_C)_j) = \sum_{n<_P j} d_j
	\]
	Since we have infinite choices of $\lambda\in \K$ but only finitely many possibilities for $\dim_\K(\lim_{n<_P j} (B(\lambda)_C)_j)$, by the \emph{Infinite Pigeonhole principle} there must be a positive integer $d_L\in\Z$, with $0\leq d_L \leq D_L$, and an infinite subset $I\subseteq\K$ such that $\dim_\K(\lim_{n<_P j} (B(\lambda_i)_C)_j) = d_L$ for any $\lambda_i \in I$. Thus $\{B(\lambda_i)_P\}_{\lambda_i\in I} \defeq \{E_{s\downarrow}(B(\lambda_i)_C)\}_{\lambda_i\in I}$ is a brick-continuous family with dimension vector $\underline{d}' \defeq (\underline{d}, d_L)$ (and dimension $d+ d_L$).
	
	If we suppose that the result holds for $m-1$ reduction steps, with an analogous reasoning we conclude by induction that $\{\widetilde{E}_\downarrow (B(\lambda)_C)\}_{\lambda\in \K}$ contains a brick continuous family $\{B(\lambda_i)_P\}_{\lambda_i\in I}$ ($I\subseteq \K$ infinite subset) with a fixed dimension vector.	
\end{proof}

We can summarise our discussion with the following theorem:

\begin{theorem}\label{thm:genbrickinfinite}
	Any representation-infinite poset $P$ is brick-continuous and admits a generic brick over $\incidencealgebra{P}$ that can be explicitly computed.
\end{theorem}

\begin{example}
	The poset $P$ can be reduced to the critical poset $\qDtilde_4$ using a contraction	
	
	% https://q.uiver.app/#q=WzAsMTQsWzAsMSwiXFxidWxsZXQiXSxbMSwxLCJcXHN0YXIiXSxbMiwwLCJcXGJ1bGxldCJdLFszLDEsIlxcc3RhciJdLFsyLDIsIlxcc3RhciJdLFsyLDMsIlxcYnVsbGV0Il0sWzQsMSwiXFxidWxsZXQiXSxbNywxLCJcXGJ1bGxldCJdLFs4LDEsIlxcc3RhciJdLFs4LDAsIlxcYnVsbGV0Il0sWzgsMiwiXFxidWxsZXQiXSxbOSwxLCJcXGJ1bGxldCJdLFs1LDEsIlAiXSxbNiwxLCJcXHFEdGlsZGVfNCJdLFswLDFdLFsxLDJdLFszLDJdLFs0LDNdLFs0LDIsIiIsMSx7InN0eWxlIjp7ImJvZHkiOnsibmFtZSI6ImRhc2hlZCJ9LCJoZWFkIjp7Im5hbWUiOiJub25lIn19fV0sWzMsNl0sWzUsNF0sWzcsOF0sWzEwLDhdLFs4LDExXSxbOCw5XSxbNCwxXSxbMTIsMTNdXQ==
	\[\begin{tikzcd}[ampersand replacement=\&,cramped,sep=small]
		\&\& \bullet \&\&\&\&\&\& \bullet \& \\
		\bullet \& \star \&\& \star \& \bullet \& P \& {\qDtilde_4} \& \bullet \& \star \& \bullet \\
		\&\& \star \&\&\&\&\&\& \bullet \\
		\&\& \bullet
		\arrow[from=2-1, to=2-2]
		\arrow[from=2-2, to=1-3]
		\arrow[from=2-4, to=1-3]
		\arrow[from=2-4, to=2-5]
		\arrow[from=2-6, to=2-7]
		\arrow[from=2-8, to=2-9]
		\arrow[from=2-9, to=1-9]
		\arrow[from=2-9, to=2-10]
		\arrow[dashed, no head, from=3-3, to=1-3]
		\arrow[from=3-3, to=2-2]
		\arrow[from=3-3, to=2-4]
		\arrow[from=3-9, to=2-9]
		\arrow[from=4-3, to=3-3]
	\end{tikzcd}\]
	
	\noindent so it is representation-infinite by Theorem \ref{thm:Lou75concealed}. Thus, we can apply Theorem \ref{thm:genbrickinfinite} to obtain a generic brick for the poset $P$:
	
	% https://q.uiver.app/#q=WzAsMTQsWzAsMSwiSyh4KSJdLFsxLDEsIksoeCleMiJdLFsxLDAsIksoeCkiXSxbMSwyLCJLKHgpIl0sWzIsMSwiSyh4KSJdLFs1LDEsIksoeCkiXSxbNiwxLCJLKHgpXjIiXSxbNywwLCJLKHgpIl0sWzgsMSwiSyh4KV4yIl0sWzcsMiwiSyh4KV4yIl0sWzcsMywiSyh4KSJdLFs5LDEsIksoeCkiXSxbMywxLCJHX3tcXHFEdGlsZGVfNH0iXSxbNCwxLCJcXHdpZGV0aWxkZXtFfShHX3tcXHFEdGlsZGVfNH0pIl0sWzAsMSwiXFxsZWZ0KFxcYmVnaW57c21hbGxtYXRyaXh9IDAgXFxcXCAxIFxcZW5ke3NtYWxsbWF0cml4fSBcXHJpZ2h0KSJdLFszLDEsIlxcbGVmdChcXGJlZ2lue3NtYWxsbWF0cml4fSAxIFxcXFwgMCBcXGVuZHtzbWFsbG1hdHJpeH0gXFxyaWdodCkiXSxbMSw0LCJcXGxlZnQoXFxiZWdpbntzbWFsbG1hdHJpeH0gMSAmIHggXFxlbmR7c21hbGxtYXRyaXh9IFxccmlnaHQpIiwyXSxbMSwyLCJcXGxlZnQoXFxiZWdpbntzbWFsbG1hdHJpeH0gMSAmIDEgXFxlbmR7c21hbGxtYXRyaXh9IFxccmlnaHQpIiwyXSxbNiw3LCJcXGxlZnQoXFxiZWdpbntzbWFsbG1hdHJpeH0gMSAmIDEgXFxlbmR7c21hbGxtYXRyaXh9IFxccmlnaHQpIl0sWzgsNywiXFxsZWZ0KFxcYmVnaW57c21hbGxtYXRyaXh9IDEgJiAxIFxcZW5ke3NtYWxsbWF0cml4fSBcXHJpZ2h0KSIsMl0sWzksOCwiMSIsMl0sWzksNywiIiwxLHsic3R5bGUiOnsiYm9keSI6eyJuYW1lIjoiZGFzaGVkIn0sImhlYWQiOnsibmFtZSI6Im5vbmUifX19XSxbMTAsOSwiXFxsZWZ0KFxcYmVnaW57c21hbGxtYXRyaXh9IDEgXFxcXCAwIFxcZW5ke3NtYWxsbWF0cml4fSBcXHJpZ2h0KSJdLFs5LDYsIjEiXSxbNSw2LCJcXGxlZnQoXFxiZWdpbntzbWFsbG1hdHJpeH0gMCBcXFxcIDEgXFxlbmR7c21hbGxtYXRyaXh9IFxccmlnaHQpIl0sWzgsMTEsIlxcbGVmdChcXGJlZ2lue3NtYWxsbWF0cml4fSAxICYgeCBcXGVuZHtzbWFsbG1hdHJpeH0gXFxyaWdodCkiXSxbMTIsMTMsIlxcd2lkZXRpbGRle0V9IiwwLHsic3R5bGUiOnsidGFpbCI6eyJuYW1lIjoibWFwcyB0byJ9fX1dXQ==
	\[\begin{tikzcd}[ampersand replacement=\&,cramped,sep=small]
		\& {K(x)} \&\&\&\&\&\& {K(x)} \&\& \\
		{K(x)} \& {K(x)^2} \& {K(x)} \& {G_{\qDtilde_4}} \& {\widetilde{E}(G_{\qDtilde_4})} \& {K(x)} \& {K(x)^2} \&\& {K(x)^2} \& {K(x)} \\
		\& {K(x)} \&\&\&\&\&\& {K(x)^2} \\
		\&\&\&\&\&\&\& {K(x)}
		\arrow["\begin{array}{c} \left(\begin{smallmatrix} 0 \\ 1 \end{smallmatrix} \right) \end{array}", from=2-1, to=2-2]
		\arrow["{\left(\begin{smallmatrix} 1 & 1 \end{smallmatrix} \right)}"', from=2-2, to=1-2]
		\arrow["{\left(\begin{smallmatrix} 1 & x \end{smallmatrix} \right)}"', from=2-2, to=2-3]
		\arrow["{\widetilde{E}}", maps to, from=2-4, to=2-5]
		\arrow["\begin{array}{c} \left(\begin{smallmatrix} 0 \\ 1 \end{smallmatrix} \right) \end{array}", from=2-6, to=2-7]
		\arrow["{\left(\begin{smallmatrix} 1 & 1 \end{smallmatrix} \right)}", from=2-7, to=1-8]
		\arrow["{\left(\begin{smallmatrix} 1 & 1 \end{smallmatrix} \right)}"', from=2-9, to=1-8]
		\arrow["{\left(\begin{smallmatrix} 1 & x \end{smallmatrix} \right)}", from=2-9, to=2-10]
		\arrow["\begin{array}{c} \left(\begin{smallmatrix} 1 \\ 0 \end{smallmatrix} \right) \end{array}", from=3-2, to=2-2]
		\arrow[dashed, no head, from=3-8, to=1-8]
		\arrow["1", from=3-8, to=2-7]
		\arrow["1"', from=3-8, to=2-9]
		\arrow["\begin{array}{c} \left(\begin{smallmatrix} 1 \\ 0 \end{smallmatrix} \right) \end{array}", from=4-8, to=3-8]
	\end{tikzcd}\]
	
\end{example}
\smallskip 

Applying ideas similar to those used previously, we can easily prove the following:

\begin{theorem}\label{thm:generalisation}
	Let $\Lambda, \Lambda'$ be finite-dimensional algebras.
	\begin{enumerate}
		\item If $\Lambda'$ is brick-continuous and $F \colon \Modf{\Lambda'} \to \Modf{\Lambda}$ is a fully faithful exact functor, then $\Lambda$ is brick-continuous.
		\item If $\Lambda'$ has a generic brick and $F \colon \Mod{\Lambda'} \to \Mod{\Lambda}$ is a fully faithful exact functor that preserves endofinite modules, then $\Lambda$ admits a generic brick.
	\end{enumerate}
\end{theorem}
\begin{proof}
	$(1)$ By Lemma \ref{lemma:brickcontdef} we can suppose to have a brick-continuous family $\{B_i\}_{i\in I}$ of $\Lambda'$ with the same dimension vector. Then $\{F(B_i)\}_{i\in I}$ is a family of non-isomorphic bricks (by fully faithfulness) with the same dimension vector (since exactness gives us a group homomorphism of the corresponding Grothendieck groups \cite[Ch.II, 6.1.5]{weibel2013k}).
	
	$(2)$ Let $G$ be a generic brick of $\Lambda'$. Then $F(G)$ is a brick (by fully faithfulness) of infinite length (by fully faithfulness and exactness) with finite endolength (since $F$ preserves endofinite modules), so a generic brick for $\Lambda$.
\end{proof}

\begin{remark}
	Theorem \ref{thm:generalisation} could be used to study properties of algebras \emph{linked} to incidence algebras:
	\begin{itemize}
		\itemsep=0em
		\item If $P$ is a representation-infinite poset and $f\colon\Lambda \to \incidencealgebra{P}$ is a ring epimorphism, then $\Lambda$ is brick-continuous and admits a generic brick. Indeed we can apply the Theorem to $F=f_*\colon \Mod{\incidencealgebra{P}} \to \Mod{\Lambda}$ (and $F=f_*\colon \Modf{\incidencealgebra{P}} \to \Modf{\Lambda}$).
		\item In general if $\Lambda'=\incidencealgebra{P}$ for a representation-infinite poset $P$, then the existence of a functor $F$ as in the Theorem gives us some information about $\Lambda$.
	\end{itemize}
\end{remark}

\begin{example}
	Consider the poset $\qR_3$. We have the ring epimorphism $f\colon \palgebra{\qR_3}\to \palgebra{\qR_3}/I_{\qR_3}=\incidencealgebra{\qR_3}$. The generic brick $G_{\qR_3}$ gives us a generic brick for the algebra $\palgebra{\qR_3}$:
	
	% https://q.uiver.app/#q=WzAsMTgsWzIsMSwiXFxLKHgpIl0sWzMsMSwiXFxLKHgpXjIiXSxbNCwxLCJcXEsoeCleMyJdLFs1LDIsIlxcSyh4KV4yIl0sWzUsMCwiXFxLKHgpXjIiXSxbNiwxLCJcXEsoeCleMyJdLFs3LDEsIlxcSyh4KV4yIl0sWzgsMSwiXFxLKHgpIl0sWzAsMSwiXFxNb2R7XFxpbmNpZGVuY2VhbGdlYnJhe1xcS317XFxxUl8zfX0iXSxbMiw0LCJcXEsoeCkiXSxbMyw0LCJcXEsoeCleMiJdLFs0LDQsIlxcSyh4KV4zIl0sWzUsNSwiXFxLKHgpXjIiXSxbNSwzLCJcXEsoeCleMiJdLFs2LDQsIlxcSyh4KV4zIl0sWzcsNCwiXFxLKHgpXjIiXSxbOCw0LCJcXEsoeCkiXSxbMCw0LCJcXE1vZHtcXHBhbGdlYnJhe1xccVJfM319Il0sWzAsMSwiXFxwaGlfezEsMX0iXSxbMSwyLCJcXHBoaV97MiwxfSJdLFszLDIsIlxccHNpX3syLDF9Il0sWzIsNCwiXFxiZXRhX3giXSxbMyw1LCJcXHBoaV97MiwxfSIsMl0sWzUsNCwiXFxhbHBoYSIsMl0sWzYsNSwiXFxwc2lfezIsMX0iLDJdLFs3LDYsIlxccHNpX3sxLDF9IiwyXSxbMyw0LCIiLDEseyJzdHlsZSI6eyJib2R5Ijp7Im5hbWUiOiJkYXNoZWQifSwiaGVhZCI6eyJuYW1lIjoibm9uZSJ9fX1dLFs5LDEwLCJcXHBoaV97MSwxfSJdLFsxMCwxMSwiXFxwaGlfezIsMX0iXSxbMTIsMTEsIlxccHNpX3syLDF9Il0sWzExLDEzLCJcXGJldGFfeCJdLFsxMiwxNCwiXFxwaGlfezIsMX0iLDJdLFsxNCwxMywiXFxhbHBoYSIsMl0sWzE1LDE0LCJcXHBzaV97MiwxfSIsMl0sWzE2LDE1LCJcXHBzaV97MSwxfSIsMl0sWzgsMTcsImZfKiIsMl1d
	\[\begin{tikzcd}[ampersand replacement=\&,cramped,sep=small]
		\&\&\&\&\& {\K(x)^2} \&\&\& \\
		{\Mod{\incidencealgebra{\qR_3}}} \&\& {\K(x)} \& {\K(x)^2} \& {\K(x)^3} \&\& {\K(x)^3} \& {\K(x)^2} \& {\K(x)} \\
		\&\&\&\&\& {\K(x)^2} \\
		\&\&\&\&\& {\K(x)^2} \\
		{\Mod{\palgebra{\qR_3}}} \&\& {\K(x)} \& {\K(x)^2} \& {\K(x)^3} \&\& {\K(x)^3} \& {\K(x)^2} \& {\K(x)} \\
		\&\&\&\&\& {\K(x)^2}
		\arrow["{f_*}"', from=2-1, to=5-1]
		\arrow["{\phi_{1,1}}", from=2-3, to=2-4]
		\arrow["{\phi_{2,1}}", from=2-4, to=2-5]
		\arrow["{\beta_x}", from=2-5, to=1-6]
		\arrow["\alpha"', from=2-7, to=1-6]
		\arrow["{\psi_{2,1}}"', from=2-8, to=2-7]
		\arrow["{\psi_{1,1}}"', from=2-9, to=2-8]
		\arrow[dashed, no head, from=3-6, to=1-6]
		\arrow["{\psi_{2,1}}", from=3-6, to=2-5]
		\arrow["{\phi_{2,1}}"', from=3-6, to=2-7]
		\arrow["{\phi_{1,1}}", from=5-3, to=5-4]
		\arrow["{\phi_{2,1}}", from=5-4, to=5-5]
		\arrow["{\beta_x}", from=5-5, to=4-6]
		\arrow["\alpha"', from=5-7, to=4-6]
		\arrow["{\psi_{2,1}}"', from=5-8, to=5-7]
		\arrow["{\psi_{1,1}}"', from=5-9, to=5-8]
		\arrow["{\psi_{2,1}}", from=6-6, to=5-5]
		\arrow["{\phi_{2,1}}"', from=6-6, to=5-7]
	\end{tikzcd}\]	
	
\end{example}

